\chapter{Binary Quadratic Forms}

\section{The Quadratic Form Associated to a Lucas Sequence}

In this section, we consider a Lucas sequence $(u_n)$ given by
\begin{equation*}
  u_0:=0,\qquad u_1:=1,\qquad\text{and}\qquad u_{n+2}:=a u_{n+1} -b u_n,
\end{equation*}
where $a$ and $b$ are integers.

\begin{definition}
  The \emph{characteristic form} of a Lucas sequence $(u_n)$ is defined by
  \begin{equation*}
    f_u(x,y):=x^2-axy+by^2.
  \end{equation*}
  The discriminant $\Delta_u=a^2-4b$ of the characteristic form $f_u(x,y)$ is called the discriminant of the Lucas sequence $(u_n)$.
\end{definition}

We saw in \cref{theorem:binets-formula}, that the entries in the sequence $(u_n)$ satisfy Binet's formula
\begin{equation*}
  u_n=\frac{\alpha^n-\beta^n}{\alpha-\beta},
\end{equation*}
where $\alpha=\frac{a+\sqrt{\Delta_u}}{2}$ and $\beta=\frac{a-\sqrt{\Delta_u}}{2}$. In the following theorem, we prove a further relation between the Lucas sequence $(u_n)$ and its characteristic form.

\begin{example}
  The Fibonacci sequence is the Lucas sequence defined by $a=1$ and $b=-1$. Its recurrence is given by
  \begin{equation*}
    F_{n+2} = 1 F_{n+1} - (-1)F_{n}.
  \end{equation*}
  Thus, the characteristic form of the Fibonacci sequence is
  \begin{equation*}
    x^2 - xy -y^2,
  \end{equation*}
  which has discriminant $\Delta=1^2 -4(-1)=5$.
\end{example}

\begin{theorem}
  Let $f_u(x,y)$ be the characteristic form of the Lucas sequence $(u_n)$. Then the entries in $(u_n)$ satisfy the equation
  \begin{equation*}
    f_u(u_{n+1},u_n)= b^n.
  \end{equation*}
\end{theorem}

\begin{proof}
  The proof is by induction on $n$. In the base case, we have
  \begin{equation*}
    f_u(u_{1},u_0)=u_{1}^2 - au_{1}u_0 + bu_0^2=1=b^0.
  \end{equation*}
  For the inductive step, assume that the identity $f_u(u_{n+1},u_n)= b^n$ holds. Then we calculate:
  \begin{align*}
    f_u(u_{n+2},u_{n+1})
    & = u_{n+2}^2 - au_{n+2}u_{n+1} + bu_{n+1}^2 \\
    & = (a u_{n+1} -b u_n)^2 - a(a u_{n+1} -b u_n)u_{n+1} + b u_{n+1}^2 \\
    & = -ab u_{n+1}u_n + b^2 u_n^2 + bu_{n+1}^2 \\
    & = b(u_{n+1}^2 - au_{n+1}u_n + b u_n^2) \\
    & = b^{n+1}.\qedhere
  \end{align*}
\end{proof}

%%% Local Variables:
%%% mode: latex
%%% TeX-master: "elementary-number-theory"
%%% TeX-engine: luatex
%%% End:
